% Created 2026-04-02 Thu 09:24
% Intended LaTeX compiler: pdflatex
\documentclass[11pt]{article}
\usepackage[utf8]{inputenc}
\usepackage[T1]{fontenc}
\usepackage{graphicx}
\usepackage{longtable}
\usepackage{wrapfig}
\usepackage{rotating}
\usepackage[normalem]{ulem}
\usepackage{amsmath}
\usepackage{amssymb}
\usepackage{capt-of}
\usepackage{hyperref}
\author{David Jácome}
\date{\today}
\title{Probabilidad y Estadística-Clase 01}
\hypersetup{
 pdfauthor={David Jácome},
 pdftitle={Probabilidad y Estadística-Clase 01},
 pdfkeywords={},
 pdfsubject={},
 pdfcreator={Emacs 29.3 (Org mode 9.6.15)}, 
 pdflang={English}}
\begin{document}

\maketitle
\tableofcontents


\section{Sílabo}
\label{sec:orge993f0e}
\subsection{Primer Parcial}
\label{sec:orga9bfa66}
\begin{enumerate}
\item Estadística descriptiva
\item Probabilidad
\end{enumerate}
\subsection{Segundo Parcial}
\label{sec:orgc8588ab}
\begin{enumerate}
\item Probabilidad (parte 2)
\item Variables aleatorias
\item Distribuciones especiales
\item Distribuciones de muestreo (TLC)
\item Intervalos de confianza
\end{enumerate}

\section{Estadística Descriptiva}
\label{sec:orgc09a1a2}
\subsection{Variable}
\label{sec:org561a64b}
\begin{itemize}
\item Puede ser cualquier letra mayúscula.
\item Es una característica de estudio.
\end{itemize}

\subsection{Conceptos Básicos}
\label{sec:orgd4cfbe8}
\begin{itemize}
\item \textbf{Población (\(N\)):} El conjunto total de datos.
\item \textbf{Muestra (\(n\)):} Un subconjunto de la población.
\item \textbf{Individuo:} Cada elemento o dato individual.
\end{itemize}

\begin{quote}
"Cada individuo es un dato."
\end{quote}

\subsection{Frecuencia}
\label{sec:orgb49ffc1}
La frecuencia se divide principalmente en dos tipos:

\begin{enumerate}
\item \textbf{Absoluta (\(F_i, n_i\)):} Consiste en contar los datos.
\begin{itemize}
\item \(n \to\) Muestra
\item \(N \to\) Población
\end{itemize}
\item \textbf{Relativa (\(F_{ri}, f_i\)):} Es una proporción o probabilidad.
$$\displaystyle f_{ri} = \frac{f_i}{n}$$
Donde \(n\) es el número total de datos.
\end{enumerate}

\subsection{Presentación de Datos}
\label{sec:org6b37bc0}
Existen tres formas principales de presentar la información:
\begin{enumerate}
\item \textbf{Individual}
\item \textbf{Agrupadas:}
\begin{itemize}
\item Simples
\item Intervalos de clase (con tablas de frecuencia)
\end{itemize}
\item \textbf{Gráfica}
\end{enumerate}

\subsection{1. Individuales}
\label{sec:orgc4e9a88}
\(x_1, x_2, x_3, x_4, x_5, x_6, x_7, \dots, x_{final}\)

\subsection{2. Tablas de Frecuencia}
\label{sec:orgc2245f1}
\(i\): Interés

\subsubsection{Tabla de agrupado simple}
\label{sec:org302b0e6}
\begin{center}
\begin{tabular}{rlll}
\(i\) & \(x_i\) & \(f_i\) & \(f_{ri}\)\\[0pt]
\hline
1 & \(x_1\) & \(f_1\) & \(f_{r1}\)\\[0pt]
2 & \(x_2\) & \(f_2\) & \(f_{r2}\)\\[0pt]
3 & \(x_3\) & \(f_3\) & \(f_{r3}\)\\[0pt]
4 & \(x_4\) & \(f_4\) & \(f_{r4}\)\\[0pt]
5 & \(x_5\) & \(f_5\) & \(f_{r5}\)\\[0pt]
\(\dots\) & \(\dots\) & \(\dots\) & \(\dots\)\\[0pt]
\(n\) & \(x_n\) & \(f_k\) & \(f_{rk}\)\\[0pt]
\end{tabular}
\end{center}

\subsubsection{Tabla por intervalos de clase}
\label{sec:orgb146d11}
\begin{center}
\begin{tabular}{rlllll}
\(i\) & \(L_i\) & \(L_{i+1} = L_s\) & \(m_i\) & \(f_i\) & \(f_{ri}\)\\[0pt]
\hline
1 & \(L_1\) & \(L_2\) & \(m_1\) & \(f_1\) & \(f_{r1}\)\\[0pt]
2 & \(L_2\) & \(L_3\) & \(m_2\) & \(f_2\) & \(f_{r2}\)\\[0pt]
3 & \(L_3\) & \(L_4\) & \(m_3\) & \(f_3\) & \(f_{r3}\)\\[0pt]
4 & \(L_4\) & \(L_5\) & \(m_4\) & \(f_4\) & \(f_{r4}\)\\[0pt]
\(\dots\) & \(\dots\) & \(\dots\) & \(\dots\) & \(\dots\) & \(\dots\)\\[0pt]
\(k\) & \(L_k\) & \(L_{k+1}\) & \(m_k\) & \(f_k\) & \(f_{rk}\)\\[0pt]
\end{tabular}
\end{center}

\begin{itemize}
\item \(L_i\): Límite inferior.
\item \(L_s\): Límite superior.
\item \textbf{Punto medio o marca de clase (\(m_i\)):} $$\displaystyle m_i = \frac{L_i + L_s}{2}$$
\end{itemize}

\section{Medidas Descriptivas}
\label{sec:org22e6d95}
\subsection{Medidas de tendencia central}
\label{sec:org95309da}
\begin{itemize}
\item \textbf{Media aritmética / Promedio / Esperanza} (\(\bar{x}\) o \(\mu\))
\item \textbf{Moda} (\(M_o\))
\item \textbf{Mediana / Percentil 50} (\(M_e\))
\end{itemize}

\subsection{Cálculo de la Media (\(\bar{x}\))}
\label{sec:org2e47109}

\subsubsection{Caso Individual}
\label{sec:orgca7114b}
$$\displaystyle \bar{x} = \frac{\sum_{i=1}^{n} x_i}{n}$$

\subsubsection{Tablas de Frecuencia}
\label{sec:org78f411d}
Para datos organizados en tablas, tenemos dos formas:

\begin{enumerate}
\item Agrupados simples
\label{sec:org3857200}
\begin{itemize}
\item \textbf{Forma 1:}
$$\displaystyle \bar{x} = \frac{\sum_{i=1}^{k} x_i \cdot f_i}{n}$$
\item \textbf{Forma 2:}
$$\displaystyle \bar{x} = \sum_{i=1}^{k} x_i \cdot f_{ri}$$
\end{itemize}

\item Intervalos de clase
\label{sec:org76aa100}
Se utiliza la marca de clase (\(m_i\)):
$$\displaystyle \bar{x} = \frac{\sum_{i=1}^{n} m_i \cdot f_i}{n} = \sum_{i=1}^{n} m_i \cdot f_{ri}$$
\end{enumerate}

\subsection{Moda (\(M_o\))}
\label{sec:org37db3ed}
\begin{itemize}
\item Es el dato que más se repite.
\item Es el dato con mayor frecuencia (absoluta o relativa).
\item \textbf{Intervalo modal:} En tablas de intervalos, es el intervalo con mayor frecuencia.
\end{itemize}

\subsection{Mediana (\(M_e\)): Cálculo de Percentil 50}
\label{sec:org1ff2e72}
\begin{IMPORTANT}
Los datos \textbf{deben estar ordenados} de menor a mayor.
\end{IMPORTANT}

\subsubsection{Caso Individual}
\label{sec:orgd76eaf5}
\begin{itemize}
\item \textbf{Si el \# de datos es par:} Es el promedio de los dos valores centrales.
\item \textbf{Si el \# de datos es impar:} Es el valor que se encuentra justo en medio (la mitad).
\end{itemize}


\section{Medidas de Dispersión}
\label{sec:org99c1bcd}
\subsection{Conceptos y Símbolos}
\label{sec:org99ae5c1}
\begin{center}
\begin{tabular}{lll}
Medida & Muestra & Población\\[0pt]
\hline
Varianza & \(s^2\) & \(\sigma^2\)\\[0pt]
Desviación Estándar & \(s\) & \(\sigma\)\\[0pt]
\end{tabular}
\end{center}

\begin{quote}
Relación fundamental:
\(\sqrt{Varianza} = Desviación\)
\((Desviación)^2 = Varianza\)
\end{quote}

\subsection{Varianza Muestral (\(s^2\))}
\label{sec:org23ab1ae}
\subsubsection{Caso Individual}
\label{sec:org41f5fe5}
$$\displaystyle s^2 = \frac{\sum_{i=1}^{n} (x_i - \bar{x})^2}{n - 1}$$

\subsubsection{Tablas de Frecuencia}
\label{sec:org1bc72ba}
\begin{itemize}
\item \textbf{Agrupados simples:}
$$\displaystyle s^2 = \frac{\sum_{i=1}^{k} (x_i - \bar{x})^2 \cdot f_i}{n - 1}$$
\item \textbf{Intervalos de clase:} (se usa la marca de clase \(m_i\))
$$\displaystyle s^2 = \frac{\sum_{i=1}^{k} (m_i - \bar{x})^2 \cdot f_i}{n - 1}$$
\end{itemize}

\subsection{Rango e Intervalo Intercuartil}
\label{sec:org3d44cf2}
\begin{itemize}
\item \textbf{Rango:}
$$\displaystyle R = X_{max} - X_{min}$$
\emph{(Los datos deben estar ordenados de menor a mayor)}
\item \textbf{Rango Intercuartil (\(RIQ\)):}
$$\displaystyle RIQ = Q_3 - Q_1$$
\begin{itemize}
\item \(Q_1\): Percentil 25 (\(P_{25}\))
\item \(Q_3\): Percentil 75 (\(P_{75}\))
\end{itemize}
\end{itemize}

\subsection{Coeficiente de Variación (\(CV\))}
\label{sec:org7c247d6}
Determina qué tan homogéneos o heterogéneos son los datos:
$$\displaystyle CV = \frac{s}{\bar{x}}$$

\begin{itemize}
\item \textbf{Interpretación:}
\begin{itemize}
\item Si \(CV > 0.2 \implies\) Datos \textbf{Heterogéneos}.
\item Si \(CV < 0.2 \implies\) Datos \textbf{Homogéneos}.
\end{itemize}
\end{itemize}
\end{document}
